\documentclass[a4paper]{report}
\usepackage[utf8]{inputenc}
\usepackage[portuguese]{babel}
\usepackage{hyperref}
\usepackage{a4wide}
\hypersetup{pdftitle={PCP - OpenMP},
pdfauthor={João Teixeira, José Ferreira},
colorlinks=true,
urlcolor=blue,
linkcolor=black}
\usepackage{subcaption}
\usepackage{listings}
\usepackage{booktabs}
\usepackage{multirow}
\usepackage{appendix}
\usepackage{tikz}
\usepackage{authblk}
\usepackage{bashful}
\usepackage{verbatim}
\usepackage{amssymb}
\usepackage{multirow}
\usepackage{mwe}
\usepackage[parfill]{parskip}
\usetikzlibrary{positioning,automata,decorations.markings}
\AfterEndEnvironment{figure}{\noindent\ignorespaces}
\AfterEndEnvironment{table}{\noindent\ignorespaces}

\begin{document}

\title{Paradigmas de Computação Paralela\\Bucket Sort com OpenMP}
\author{João Teixeira (A85504) \and José Filipe Ferreira (A83683)}
\date{\today}

\begin{center}
    \begin{minipage}{0.75\linewidth}
        \centering
        \includegraphics[width=0.4\textwidth]{images/eng.jpeg}\par\vspace{1cm}
        \vspace{1.5cm}
        \href{https://www.uminho.pt/PT}
        {\color{black}{\scshape\LARGE Universidade do Minho}} \par
        \vspace{1cm}
        \href{https://www.di.uminho.pt/}
        {\color{black}{\scshape\Large Departamento de Informática}} \par
        \vspace{1.5cm}
        \maketitle
    \end{minipage}
\end{center}

\tableofcontents

\pagebreak

\chapter{Introdução}
O algoritmo escolhido para o projeto da unidade curricular de Computação
Paralela e Distribuída foi o \textit{Bucket Sort}.

Inicialmente desenvolvemos uma versão sequencial do projeto e procedemos ao
\textit{benchmarking} do programa resultante. Em seguida convertemos a
implementação sequencial numa versão com utilização de memoria partilhada
fazendo uso de \textit{OpenMP} e comparamos o resultado com a versão sequencial
desenvolvida anteriormente.

Ao longo deste relatório iremos descrever a metodologia utilizada e os
resultados de \textit{benchmarking} obtidos ao longo deste projeto.

De notar que todos os benchmarks descritos foram efeutados em nós do tipo 652 do
cluster \textit{SeARCH}, e todos os executáveis foram compilados com a versão
7.2.0 do \textit{gcc}, com as flags \textit{-O3 -ftree-vectorize -fopenmp
-std=c11} para garantir a maior consistência entre os diferentes benchmarks.

\chapter{Sequencial}
O \textit{Bucket Sort} consiste em definir um conjunto de N "baldes"
inicialmente vazios. Em seguida os elementos do vetor a ser ordenado são
distribuídos pelos baldes. O critério escolhido para esta distribuição foi
calcular o máximo e o mínimo do vetor a ser ordenado e dividir os intervalos de
valores de cada balde em intervalos do mesmo tamanho. Os elementos do vetor
inicial são então distribuídos pelo respetivo balde com base no intervalo
definido. Seguidamente o conteúdo de cada balde é ordenado recorrendo ao
\textit{quicksort} presente na \textit{standard library} de C. Finalmente todos
os elementos são copiados um a um para o vetor original.

Para testar se a primeira implementação sequencial produzia de facto vetores
ordenados criamos um \textit{script} que permite gerar N testes aleatórios e
comparar o resultado do nosso programa com o resultado de ordenar os valores
com o comando \textit{sort} de \textit{bash}.

Numa implementação inicial desta versão sequencial, o algoritmo consistia em
iterar sobre o array dos elementos a ordenar e inserir o elemento a processar no
respetivo balde. Durante a paralelização deste, identificamos alguns problemas,
que serão descritos em maior pormenor no capitulo \ref{chap:omp}, que nos levou
a tomar uma abordagem diferente, em vez de iterar pelo array uma vez e colocar
os elementos no respetivo balde, cada balde irá iterar pelo array, escolhendo os
elementos que a ele lhe pertencem.

Fazer isto aumentou a complexidade do algoritmo, passando de iterar uma vez pelo
array, para N vezes, onde N é o número de baldes existentes. Como
consequência disto também os tempos de execução aumentarem.

\chapter{OpenMP} \label{chap:omp}
Fazendo uso da implementação sequencial descrita no capitulo anterior, começamos
a conversão para uma versão com um modelo de memoria partilhada fazendo uso de
\textit{OpenMP}.

Para a paralelização da primeira versão sequencial, começamos por tratar do
ciclo mais longo que itera por todos os elementos do \textit{array}. Ao fazer
uma analise mais detalhada deste, deparamo-nos com a existência de uma zona
critica onde era escrito cada elemento no respetivo balde. Para tentar resolver
o problema tentamos duas abordagens distintas, a utilização de tarefas e a
declaração da zona critica. Em ambas as opções obtivemos resultados piores
quando comparados com a sua versão sequencial, pois tanto a gestão de zonas
criticas, bem como a gestão das tarefas são tarefas pesadas.

Vendo estes resultados, tentamos adaptar o algoritmo para um formato mais
paralelizável, removendo a zona critica.

Nesta nova versão verificamos que os tempos de execução paralelos estavam a
baixar, obtendo speedups consideráveis até às 16 \textit{threads}. Embora o
speedup esteja ligeiramente abaixo do ideal e se afastar do valor quanto maior o
numero de \textit{threads}. Isto deve-se ao número de elementos por balde não
ser uniforme e o tempo de execução só conseguir ser tão baixo quanto o balde
mais lento de processar. Por isso, quanto maior o numero de \textit{threads} em
trabalho mais se notará esse desequelibrio.

Para tentar mitigar este problema tentamos utilizar o colapsamento dos dois
ciclos aninhados, o que percorre os baldes e o interior que percorre o array dos
elementos. Para manter a escrita no balde dos elementos em paralelo a
atualização do índice de escrita no balde era feita de forma atómica. Nesta
implementação, a escalabilidade do algoritmo provou-se bastante inferior,
começando a diminuir a partir das 4 threads, mantendo a tendência até às 32
threads.

Para tentar tornar o algoritmo ainda mais escalavel, tentamos variar o número de
baldes que são utilizados para distribuir os elementos a ordenar. Numa primeira
abordagem testamos com um mesmo número de baldes para vários números de threads.
Não encontramos nenhum valor que fosse o melhor para todos o número de threads,
mas vimos que quantas mais threads utilizadas, mais benéfico era o maior número
de baldes. 

Perante esta relação tentamos testar com um valor de baldes dependente do número
de threads utilizadas. Aqui vimos que, regra geral, utilizar o dobro do número
de threads garantia o maior resultado, mas tendo em conta as caracteristicas do
nosso algoritmo, nomeadamente, o facto de por cada balde é feita uma passagem
pelo array de elementos a ordenar, a partir das 32 threads o aumento da carga de
trabalho torna-se tal que o programa começa a ficar mais lento.

Juntando a informação obtida em ambos os testes, implementamos um meio termo
entre ambas as abordagens, utilizando o dobro do número de threads para o número
de baldes, e colocando um teto máximo de 32 baldes, valores para o qual obtivemos os melhores
resultados médios até as 40 threads.

\appendix

\chapter{Tempos de execução, para 100000000 elementos aleatórios, entre -1000 e 50000}
\begin{table}[h]
\centering
\begin{tabular}{|l|l|l|l|l|}
\hline
\multicolumn{5}{|c|}{Times}                                      \\ \hline
Threads & Old Critical & Old Task    & New       & New Collapsed \\ \hline
Seq     & 16,209999    & 16,209999   & 22,299999 & 22,299999     \\ \hline
2       & 43,23        & 305,169983  & 11,403177 & 11,968859     \\ \hline
4       & 100,150002   & 937,929993  & 6,880816  & 6,872915      \\ \hline
8       & 187,929993   & 2227,460205 & 4,622566  & 5,965684      \\ \hline
16      & 479,750031   & 5547,610352 & 2,362948  & 4,678829      \\ \hline
32      & 1213,959961  & 12624,78906 & 2,363651  & 4,859851      \\ \hline
64      & 1726,070068  & 11589,56934 & 2,449052  & 4,059251      \\ \hline
\end{tabular}
\caption{\label{tab:Times}Tempos de execução iniciais, com 10 baldes}
\end{table}

\begin{table}[h]
\centering
\begin{tabular}{|l|l|l|l|l|l|}
\hline
Threads & 5 buckets & 10 buckets & 20 buckets & 40 buckets & 80 bucket \\ \hline
2       & 12.520761 & 11.218051  & 12.921473  & 17.016363  & 25.398283 \\ \hline
4       & 8.365767  & 6.739361   & 6.511606   & 8.512317   & 12.743460 \\ \hline
8       & 5.218658  & 4.984042   & 3.952061   & 4.323404   & 6.438317  \\ \hline
16      & 5.344770  & 3.167999   & 3.042659   & 2.632886   & 3.465635  \\ \hline
32      & 5.248249  & 3.147622   & 2.148584   & 2.392601   & 2.836905  \\ \hline
64      & 5.270648  & 3.161402   & 2.366316   & 2.002006   & 3.260201  \\ \hline
\end{tabular}
\caption{\label{tab:fix_size}Tempos de execução, fazendo variar o
    número de baldes utilizados}
\end{table}

\begin{table}[h]
\centering
\begin{tabular}{|l|l|l|l|l|l|}
\hline
Threads & n\_threads & 2*n\_threads & 4*n\_threads & 8*n\_threads & 16*n\_threads \\ \hline
2       & 11.395469  & 10.310162    & 10.808166    & 12.206050    & 15.356203     \\ \hline
4       & 6.129920   & 5.490817     & 6.108120     & 7.681805     & 11.012663     \\ \hline
8       & 3.719489   & 3.100843     & 3.883367     & 5.560686     & 9.027642      \\ \hline
16      & 2.217584   & 2.001313     & 2.854720     & 4.599587     & 8.157944      \\ \hline
32      & 1.824530   & 2.292776     & 3.511230     & 6.276313     & 11.743461     \\ \hline
64      & 2.486487   & 3.223805     & 5.932982     & 10.554493    & 20.222078     \\ \hline
\end{tabular}
\caption{\label{tab:var_size}Tempos de execução, fazendo variar o
    número de baldes, com base no número de threads.}
\end{table}

\begin{table}[h]
\centering
\begin{tabular}{|l|l|l|l|l|l|l|}
\hline
Threads & 30       & 32       & 34       & 36       & 38       & 40       \\ \hline
16      & 2.085598 & 2.008397 &          &          &          &          \\ \hline
32      & 1.846964 & 1.813501 & 2.572607 & 2.442342 & 2.408187 & 2.349852 \\ \hline
40      & 1.907169 & 1.799508 & 1.779662 & 1.718217 & 1.788759 & 1.881975 \\ \hline
64      & 2.599405 & 2.536593 & 2.106488 & 1.911450 & 1.976940 & 1.977957 \\ \hline
\end{tabular}
\caption{\label{tab:var_size}Tempos de execução, com o número de baldes a ser o
    dobro do número de threads, e variando o número máximo de baldes
    utilizados.}
\end{table}
\end{document}
